\section{Localizing Goal Heads}
\label{app:head_loc}

We locate the \goalheads with a simpler contrastive setup where we ask the LM to implement two different tasks/predicates $g_\mathrm{src}$ and $g_\mathrm{dst}$ in two different programming languages $\mathrm{lang}_\mathrm{src}$ and $\mathrm{lang}_\mathrm{dst}$. We aim to locate the representations that mediate the task information. Attention heads are a natural candidate as they are the only way a transfer LM can pool information from the context to the current token position\footnote{this assumption is not true for the newer generation of LMs with DeltaNet and our method is not directly applicable there.}. Concretely, we want to locate attention heads outputs in the source run $M(p_\mathrm{src})$, when patched to the destination run $M(p_\mathrm{dst})$, will trigger the execution of the source task $g_\mathrm{src}$ in the destination language $\mathrm{lang}_\mathrm{dst}$. We use different programming languages $(\mathrm{lang}_\mathrm{src} \neq \mathrm{lang}_\mathrm{dst})$ to separate the encoding of $g_\mathrm{src}$ from the surface level tokens that implement it. \Cref{fig:localization} shows a pair of example contrastive prompts.

In both source and destination runs the LM will need multiple tokens to implement $g_\mathrm{src}$ and $g_\mathrm{dst}$. To operationalize the idea in a single forward pass we use a third prompt that asks the LM to implement the source task $g_\mathrm{src}$ in destination language $\mathrm{lang}_\mathrm{dst}$. In the destination run and the target run we prefill in the implementation with a common prefix such that a successful implementation of $g_\mathrm{dst}$ in $M(p_\mathrm{dst})$ and $g_\mathrm{src}$ in $M(p_\mathrm{targ})$ will generate different next tokens $\mathrm{tok}_\mathrm{dst} \neq \mathrm{tok}_\mathrm{targ}$. While paring up contrastive tasks we choose them such that a potential implementation of $g_\mathrm{src}$ and $g_\mathrm{dst}$ in $\mathrm{lang}_\mathrm{dst}$ may share this common prefill. 

\begin{figure}[ht]
\centering
\includegraphics[width=\linewidth]{Figures/goalscope/head_localization.png}
\caption{Locating the goal heads. \textbf{Left:} source and destination prompts request different predicates in different programming languages. A separate target prompt requests the source predicate in the destination language. The two destination-language implementations share a common {\color{blue}prefix} and differ in the next token. We patch source head contributions into the destination using the mask-weighted intervention in \Cref{eqn:opt-int}, and jointly train the masks to promote the target continuation. \textbf{Right:} the selected 122 heads in Qwen2.5-Coder-32B.}
\label{fig:localization}
\end{figure}

To localize the heads we learn a joint $\mathrm{mask}$ over the heads, $\mathrm{mask}^{\ell j}$ dictates what portion of head $[\ell, j]$'s contribution in the source run are we bringing over in the patched run. Concretely we perform the following intervention on the patched run for each attention head:
\begin{equation}
a_{-1}^{\ell j}\leftarrow \mathrm{mask}^{\ell j}a_{\mathrm{src}}^{\ell j}+(1-\mathrm{mask}^{\ell j})a_{-1}^{\ell j}.
\label{eqn:opt-int-app}
\end{equation}

Here, the vector $a_{-1}^{\ell j}$ denotes the OV output of head $[\ell, j]$, at the last token position and $a_\mathrm{src}^{\ell j}$ is the same head's output in the source run, also at the last token position. The terms $\mathrm{mask}^{\ell j}$ are binary values, jointly learned with an objective to match the KL divergence over the next token distributions produced in the target run. Fromally, the loss is:

\begin{equation}
\mathcal{L}=\frac{1}{|\mathcal{B}|}\sum_{x\in\mathcal{B}}
D_{\mathrm{KL}}\!\bigg(M\left(p_\mathrm{dest}\right)\left[\leftarrow a_\mathrm{src}\right]\,\|\, M\left(p_\mathrm{targ}\right) \bigg)
+\lambda_{\mathrm{sparse}}\frac{1}{|\mathcal{H}|}\sum_{[\ell, j]\in\mathcal{H}}\mathrm{mask}^{\ell j}.
\end{equation}

We can tune the sparsity loss $\lambda_\mathrm{sparse}$ to control the number of heads being selected. $\mathcal{H}$ contains all candidate heads and $\mathcal{B}$ is a training batch. We clip $\mathrm{mask}^{\ell j}$ to stay in the $[0, 1]$ range after each optimization step.
 
% We localize heads through a joint intervention on contrastive single-predicate prompts. The source and destination request different predicates in different programming languages. A third, target prompt requests the source predicate in the destination language. The destination and target implementations share a prefix but disagree at the next token. We select contrasts whose target token also differs from the source continuation, so copying the source token alone cannot solve the localization task.


% For each attention head $h$, a learned mask $m_h\in[0,1]$ interpolates between destination and source activations before the output projection. We freeze the language model and minimize

% Here $\mathcal{H}$ contains all candidate heads and $\mathcal{B}$ is a training batch. The target is the full next-token distribution of the target prompt. The sparsity term is the mean mask value, not an unnormalized sum. Optimizing the mask jointly allows several heads to contribute to the same counterfactual continuation; it does not require every selected head to have a large effect when patched alone.

% Masks start at one and are clamped to $[0,1]$ after each AdamW update. All five saved runs use learning rate $.01$ and $\lambda_{\mathrm{sparse}}=20$. The cache is split in its stored order: the final $\lfloor .2N\rfloor$ instances form validation and the remainder form training. This is an instance split, not a held-out-predicate split. It is separate from the 253-specification readout corpus.

% \begin{table}[ht]\centering\small
% \begin{tabular}{lrrrrc}\toprule
Model & Train & Validation & Epochs & Batch & Augment\\\midrule
Qwen Coder 32B & 1139 & 284 & 50 & 64 & Yes\\
Qwen Coder 14B & 997 & 249 & 50 & 32 & No\\
Qwen Instruct 32B & 945 & 236 & 100 & 16 & No\\
OLMo Instruct 32B & 1139 & 284 & 50 & 64 & Yes\\
OLMo SFT 32B & 681 & 170 & 100 & 16 & No\\
\bottomrule\end{tabular}
% \caption{Localization settings from the saved mask checkpoints. Augmented runs resample surface forms for training and validation each epoch; unaugmented runs use precomputed source activations.}
% \label{tab:localization-training}
% \end{table}

% Augmented runs use seed 123456 and choose between changing the source language (55\%), renaming variables (35\%), and toggling type hints (10\%). The training objective concerns one continuation position; subsequent GoalScope evaluation generates and judges a multi-token readout.

% \begin{table}[ht]\centering\small
% \begin{tabular}{lrrrc}\toprule
Model & Threshold & Export layers & Heads & Selection SHA-256\\\midrule
Qwen Coder 32B & .9 & 25--55 & 122 & \texttt{6e8ab55c29}\\
Qwen Coder 14B & .9 & 0--47 & 88 & \texttt{a3b383df7e}\\
Qwen Instruct 32B & .9 & 25--55 & 108 & \texttt{10a2bad540}\\
OLMo Instruct 32B & 1.0 & 10--35 & 93 & \texttt{4534b586c1}\\
OLMo SFT 32B & .6 & 10--35 & 91 & \texttt{69aa032562}\\
\bottomrule\end{tabular}
% \caption{Exported head selections used in the five-model evaluation. Layer endpoints are inclusive. Masks are thresholded within the export band; all listed memberships match their saved training checkpoints. SHA-256 prefixes identify the exact selection files. The OLMo instruction-model threshold includes values equal to one.}
% \label{tab:localization-selection}
% \end{table}

% The core Qwen2.5-Coder-32B selection contains 122 heads, with selected layers spanning 28--55. The broader ablation set takes the top 250 mask values across all layers. Its depth coverage differs from the core selection. The small-head-budget sweep uses available snapshots and mask selections; these sets are not all nested prefixes of a single ranking.

